% !PRINTED 78
\weberpage{78}
%<<CONT

\begin{equation}
        a_{\alpha_1}^{(\beta_1)}a_{\alpha_2}^{(\beta_2)}\cdots a_{\alpha_\nu}^{(\beta_\nu)}\,
        A_{\alpha_1,\alpha_2,\ldots,\alpha_\nu}^{\beta_1,\beta_2,\ldots,\beta_\nu}
\tag{2}
\end{equation}
the aggregate of those terms of the determinant which contain the product of these elements as a factor.

By interchanging rows and columns, whereby at most the sign of the determinant is changed, one can always cause the elements
\begin{equation}
        a_{\alpha_1}^{(\beta_1)},
        a_{\alpha_2}^{(\beta_2)},\ldots,
        a_{\alpha_\nu}^{(\beta_\nu)}
\tag{3}
\end{equation}
to take the place of the elements. % [sic] Weber prints a full stop in mid-sentence before the following display; retained source-faithfully
\[
        a_1^{(1)},a_2^{(2)},\ldots,a_\nu^{(\nu)}
\]
Then rule I applies to the determination of \(A_{\alpha_1,\alpha_2,\ldots,\alpha_\nu}^{\beta_1,\beta_2,\ldots,\beta_\nu}\), and one obtains:

\gesperrt{II. Apart from sign, \(A_{\alpha_1,\alpha_2,\ldots,\alpha_\nu}^{\beta_1,\beta_2,\ldots,\beta_\nu}\) is obtained as an \((n-\nu)\)-rowed determinant by deleting from \(A\) all rows and all columns which meet in one of the elements (3), and by leaving the remaining rows and columns in their order.}

For the determination of the sign one uses the following rule.  Arrange the lower and upper indices \(1,2,\ldots,n\) in the orders
\begin{equation}
        \alpha_1,\alpha_2,\ldots,\alpha_\nu,\alpha_{\nu+1},\ldots,\alpha_n,
\tag{4}
\end{equation}
\begin{equation}
        \beta_1,\beta_2,\ldots,\beta_\nu,\beta_{\nu+1},\ldots,\beta_n,
\tag{5}
\end{equation}
where the missing \(\alpha_{\nu+1},\ldots,\alpha_n\), and likewise \(\beta_{\nu+1},\ldots,\beta_n\), are taken \gesperrt{in increasing order}.

\gesperrt{III. The \((n-\nu)\)-rowed determinant described in II receives the positive or negative sign according as the two arrangements (4), (5) of the digits \(1,2,\ldots,n\) both belong to the same kind or to different kinds.}

Indeed, the determinant changes its sign through each interchange of two lower or two upper indices.  In order to reduce the general case (2) to the special case (1), one has to make just so many transpositions of upper and lower indices that the permutations (4) and (5) both pass into the original arrangement \(1,2,3,\ldots,n\); the same number of changes of sign has then taken place.
